\section{Related Work}
\label{sec:related_work}

\subsection{Hand Gesture Recognition}

Hand gesture recognition has been extensively studied using both
vision-based and sensor-based approaches. Early methods relied on colored
gloves or depth sensors~\cite{rautaray2015vision}, while modern
approaches leverage deep learning on RGB camera input.
MediaPipe~\cite{zhang2020mediapipe} provides a lightweight hand
landmark detection model that extracts 21 3D keypoints in real-time,
forming the basis for many downstream gesture classification systems.

For temporal gesture classification, Recurrent Neural Networks (RNNs)
and Long Short-Term Memory (LSTM) networks have been widely
used~\cite{pigou2018gesture}. However, Temporal Convolutional Networks
(TCNs)~\cite{bai2018empirical, lea2017temporal} have emerged as a
competitive alternative, offering parallelizable training, stable
gradient flow via residual connections, and flexible receptive fields
through dilated causal convolutions. Lea et al.~\cite{lea2017temporal}
demonstrated that TCNs outperform LSTMs on action segmentation tasks
while being significantly faster to train.

\subsection{Model Compression for Edge Deployment}

Deploying deep learning models on resource-constrained devices (e.g.,
smartphones, embedded systems) requires model compression techniques.
Structured pruning~\cite{li2017pruning} removes entire convolutional
filters, producing genuinely smaller models without requiring sparse
hardware support. Network quantization~\cite{jacob2018quantization}
reduces numerical precision from 32-bit floating point to 8-bit
integers, achieving approximately $4\times$ memory reduction. The ONNX
Runtime~\cite{onnxruntime} provides a portable inference engine that
supports both pruned and quantized models across diverse hardware
platforms.

\subsection{Contactless Device Interaction}

Gesture-based device interaction has been explored in various contexts.
Gesture control for smart TVs~\cite{samsung2013gesture}, automotive
interfaces~\cite{ohn2014hand}, and augmented reality
headsets~\cite{meta2023hand} demonstrates the practical applicability
of hand gesture recognition. However, using gestures specifically for
\emph{cross-device content transfer} remains relatively unexplored.
Existing wireless sharing solutions such as Apple AirDrop, Google Nearby
Share, and Huawei Share rely on explicit user actions (button taps,
device selection) rather than natural gestures. GrabDrop fills this gap
by combining real-time gesture recognition with a zero-configuration
network protocol.
